\chapter{Dot Products, Length, and Angle}

So far our vectors have lived a slightly lonely life. We have added them, scaled
them, and stacked them into linear combinations -- but we have never once asked a
question of the form ``how \emph{long} is this arrow?'' or ``what is the
\emph{angle} between these two arrows?'' Those are geometric questions, and up to
now our algebra has had no way to answer them. Addition and scaling alone simply
do not know about length or angle.

This chapter introduces the one new operation that answers both questions: the
\textbf{dot product}. It is a humble-looking gadget -- multiply matching
coordinates and add them up -- but it carries the entire metric geometry
of $\R^n$ inside it. Length is a dot product. Angle is a dot product. The word
``perpendicular'' becomes the single equation ``dot product equals zero.'' And
once you see that the same little number gives the cosine of an angle, you have
in your hands the machinery behind projections, least squares, Fourier series, and the
``cosine similarity'' a search engine uses to decide that two documents are
about the same thing.

We will build it geometrically first -- arrows, lengths, the law of cosines --
then make it precise, then push it to the two great inequalities (Cauchy--Schwarz
and the triangle inequality) that hold in \emph{every} dimension. Slow down here.
The dot product is small, but it is the hinge the rest of the book swings on.

\chapterline

\section{The dot product, two ways}

\subsection{The algebraic definition}

Here is the whole operation in one line. Given two vectors in $\R^n$ with the same
number of entries,
\[
  \vec{u} = \colvec{u_1 \\ u_2 \\ \vdots \\ u_n}, \qquad
  \vec{v} = \colvec{v_1 \\ v_2 \\ \vdots \\ v_n},
\]
their \textbf{dot product} is the single number
\[
  \vec{u}\cdot\vec{v} \;=\; u_1 v_1 + u_2 v_2 + \cdots + u_n v_n
  \;=\; \sum_{i=1}^{n} u_i v_i.
\]
Multiply matching entries, add up the results. The inputs are two vectors; the
output is one \emph{scalar}. For that reason the dot product is sometimes called
the \emph{scalar product}. (You will also see it written $\inner{\vec u}{\vec v}$,
the ``inner product'' notation, and -- once we meet matrices -- as $\vec u\T\vec v$.
All three mean the same arithmetic.)

\begin{example}\label{ch04:first-dot}
Let $\vec u = \colvec{3 \\ 4}$ and $\vec v = \colvec{2 \\ -1}$. Then
\[
  \vec u \cdot \vec v = (3)(2) + (4)(-1) = 6 - 4 = 2.
\]
And in three dimensions, with $\vec a = \colvec{1\\2\\2}$ and
$\vec b = \colvec{2\\0\\-1}$,
\[
  \vec a\cdot\vec b = (1)(2) + (2)(0) + (2)(-1) = 2 + 0 - 2 = 0.
\]
That zero is not an accident; by the end of this chapter you will read it as
``$\vec a$ and $\vec b$ are perpendicular.''
\end{example}

\begin{focusbox}
The dot product eats two vectors of the \emph{same length} and returns one number.
You can only dot $\vec u\cdot\vec v$ when $\vec u$ and $\vec v$ have the same
number of components. Dotting a vector in $\R^2$ with a vector in $\R^3$ is
meaningless -- there is no rule that pairs up their entries.
\end{focusbox}

A few properties fall straight out of the definition, and we will lean on all of
them constantly.

\begin{formalbox}
For all vectors $\vec u, \vec v, \vec w \in \R^n$ and every scalar $c\in\R$:
\begin{enumerate}
  \item \textbf{Symmetry:} $\vec u\cdot\vec v = \vec v\cdot\vec u$.
  \item \textbf{Linearity in each slot:}
        $(\vec u + \vec w)\cdot\vec v = \vec u\cdot\vec v + \vec w\cdot\vec v$
        and $(c\,\vec u)\cdot\vec v = c\,(\vec u\cdot\vec v)$.
  \item \textbf{Positivity:} $\vec u\cdot\vec u = u_1^2 + \cdots + u_n^2 \ge 0$,
        with equality if and only if $\vec u = \vec 0$.
\end{enumerate}
\end{formalbox}

\begin{proof}
Each is one line of algebra. Symmetry: $u_iv_i = v_iu_i$ term by term.
Linearity: $(u_i+w_i)v_i = u_iv_i + w_iv_i$, then sum over $i$; and
$(cu_i)v_i = c(u_iv_i)$. Positivity: $\vec u\cdot\vec u = \sum_i u_i^2$ is a sum of
squares, hence $\ge 0$, and a sum of squares is zero exactly when every term is
zero, i.e. every $u_i=0$.
\end{proof}

Property 3 is the quiet hero of the whole subject. A sum of squares can never be
negative, and it is zero only for the zero vector. That single fact is what makes
$\vec u\cdot\vec u$ deserve to be called ``length squared,'' which is where we go
next.

\subsection{Length lives inside the dot product}

Look again at $\vec u\cdot\vec u = u_1^2 + \cdots + u_n^2$. In the plane, with
$\vec u = \colvec{u_1\\u_2}$, this is exactly $u_1^2 + u_2^2$ -- the inside of the
Pythagorean theorem. The hypotenuse of the right triangle with legs $u_1$ and
$u_2$ is the length of the arrow $\vec u$. So:
\[
  \norm{\vec u} \;=\; \sqrt{u_1^2 + u_2^2} \;=\; \sqrt{\vec u\cdot\vec u}.
\]
Nothing about that derivation cared that we were in two dimensions. We simply
\emph{define} length in $\R^n$ to be the same formula, and the definition is
forced on us the moment we want Pythagoras to keep holding.

\begin{figure}[H]
\centering
\begin{tikzpicture}[scale=1.1,>=Latex]
  \draw[->,gray!50] (-0.5,0) -- (5,0) node[right]{$x_1$};
  \draw[->,gray!50] (0,-0.5) -- (0,3.5) node[above]{$x_2$};
  \draw[->,very thick,vecblue] (0,0) -- (4,3)
        node[midway,above left]{$\vec u$};
  \draw[dashed,gray] (4,0) -- (4,3);
  \draw[<->,gray] (0,-0.35) -- (4,-0.35) node[midway,below]{$u_1=4$};
  \draw[<->,gray] (4.35,0) -- (4.35,3) node[midway,right]{$u_2=3$};
  \node[vecblue] at (1.7,1.9) {$\norm{\vec u}=5$};
  \draw[fill=black] (0,0) circle (1.2pt);
\end{tikzpicture}
\caption{Length is Pythagoras: $\norm{\vec u}=\sqrt{4^2+3^2}=\sqrt{25}=5$.
The dot product $\vec u\cdot\vec u = 16+9 = 25$ is the length \emph{squared}.}
\label{ch04:pythagoras}
\end{figure}

\begin{definition}\label{ch04:norm}
The \textbf{length} (or \textbf{norm}, or \textbf{magnitude}) of
$\vec u\in\R^n$ is
\[
  \norm{\vec u} \;=\; \sqrt{\vec u\cdot\vec u}
  \;=\; \sqrt{u_1^2 + u_2^2 + \cdots + u_n^2}.
\]
A vector with $\norm{\vec u}=1$ is called a \textbf{unit vector}.
\end{definition}

Because $\vec u\cdot\vec u \ge 0$, that square root is always of a nonnegative
number, so $\norm{\vec u}$ is a well-defined real number $\ge 0$, and it equals $0$
only for $\vec u = \vec 0$. Length behaves exactly as length should.

\begin{example}
$\norm{\colvec{3\\4}} = \sqrt{9+16} = 5$. \quad
$\norm{\colvec{1\\2\\2}} = \sqrt{1+4+4} = 3$. \quad
$\norm{\colvec{1\\1\\1\\1}} = \sqrt{4} = 2$.
The last one is a useful sanity check: the all-ones vector in $\R^n$ has length
$\sqrt n$, which grows without bound as the dimension does.
\end{example}

\begin{trapbox}
The dot product $\vec u\cdot\vec u$ is length \emph{squared}, not length. People
constantly drop the square root, or, worse, take the square root of the
\emph{individual} coordinates. There is exactly one square root, and it goes
\emph{around the whole sum}: $\norm{\vec u}=\sqrt{u_1^2+\cdots+u_n^2}$, never
$\abs{u_1}+\cdots+\abs{u_n}$. (That second thing is a different, also useful, way
to measure size -- but it is not the length we get from the dot product.)
\end{trapbox}

One small but heavily used fact: scaling a vector scales its length by the
\emph{absolute value} of the scalar. For any $c\in\R$,
\[
  \norm{c\,\vec u} = \sqrt{(c\vec u)\cdot(c\vec u)}
                   = \sqrt{c^2\,(\vec u\cdot\vec u)}
                   = \abs{c}\,\norm{\vec u}.
\]
Doubling an arrow doubles its length; negating it ($c=-1$) flips its direction but
leaves its length unchanged.

\section{The geometric meaning: length times length times cosine}

We now have an \emph{algebraic} number, $\vec u\cdot\vec v$, and we have used the
special case $\vec u\cdot\vec u$ to recover length. But what does the general dot
product of two \emph{different} vectors mean geometrically? The answer is one of
the most useful formulas in mathematics:
\[
  \boxed{\;\vec u\cdot\vec v \;=\; \norm{\vec u}\,\norm{\vec v}\,\cos\theta\;}
\]
where $\theta$ is the angle between the two arrows. Here is why it holds: the
``why'' is just the law of cosines wearing new clothes.

\begin{intuitionbox}
The dot product measures \emph{how much two arrows point the same way}, weighted by
how long they are. If they point in nearly the same direction, $\cos\theta$ is near
$+1$ and the dot product is large and positive. If they are perpendicular,
$\cos\theta = 0$ and the dot product is $0$. If they point in opposite directions,
$\cos\theta$ is near $-1$ and the dot product is large and negative. The sign of
$\vec u\cdot\vec v$ alone tells you whether the angle between the arrows is acute,
right, or obtuse (Figure~\ref{ch04:sign}).
\end{intuitionbox}

\begin{figure}[H]
\centering
\begin{tikzpicture}[scale=0.95,>=Latex]
  % three small panels: acute, right, obtuse
  \begin{scope}[xshift=0cm]
    \draw[->,gray!50] (-0.3,0) -- (2.2,0);
    \draw[->,gray!50] (0,-0.3) -- (0,2.2);
    \draw[->,very thick,vecblue] (0,0) -- (1.9,0.5);
    \draw[->,very thick,vecred]  (0,0) -- (1.3,1.5);
    \node[align=center] at (1.0,-0.8)
      {acute\\$\vec u\cdot\vec v>0$};
  \end{scope}
  \begin{scope}[xshift=4cm]
    \draw[->,gray!50] (-0.3,0) -- (2.2,0);
    \draw[->,gray!50] (0,-0.3) -- (0,2.2);
    \draw[->,very thick,vecblue] (0,0) -- (1.9,0.0);
    \draw[->,very thick,vecred]  (0,0) -- (0.0,1.9);
    \draw[gray] (0.3,0) -- (0.3,0.3) -- (0,0.3);
    \node[align=center] at (1.0,-0.8)
      {right\\$\vec u\cdot\vec v=0$};
  \end{scope}
  \begin{scope}[xshift=8cm]
    \draw[->,gray!50] (-1.4,0) -- (2.2,0);
    \draw[->,gray!50] (0,-0.3) -- (0,2.2);
    \draw[->,very thick,vecblue] (0,0) -- (1.9,0.5);
    \draw[->,very thick,vecred]  (0,0) -- (-1.1,1.5);
    \node[align=center] at (0.4,-0.8)
      {obtuse\\$\vec u\cdot\vec v<0$};
  \end{scope}
\end{tikzpicture}
\caption{The sign of the dot product is the sign of $\cos\theta$. Arrows that
lean the same way give a positive dot product (acute angle); perpendicular arrows
give zero (right angle); arrows that lean apart give a negative dot product
(obtuse angle).}
\label{ch04:sign}
\end{figure}

\subsection{Deriving the cosine formula from the law of cosines}

Place both arrows tail-to-tail at the origin. The vector $\vec u - \vec v$ is the
arrow that runs from the tip of $\vec v$ to the tip of $\vec u$ -- it closes the
triangle whose other two sides are $\vec u$ and $\vec v$ (Figure~\ref{ch04:loc}).

\begin{figure}[H]
\centering
\begin{tikzpicture}[scale=1.3,>=Latex]
  \coordinate (O) at (0,0);
  \coordinate (U) at (3,1.6);
  \coordinate (V) at (3.2,-0.4);
  \draw[->,very thick,vecblue] (O) -- (U) node[above left]{$\vec u$};
  \draw[->,very thick,vecred]  (O) -- (V) node[below right]{$\vec v$};
  \draw[->,thick,structurepurple] (V) -- (U)
        node[midway,right]{$\vec u-\vec v$};
  \pic[draw,->,"$\theta$",angle radius=0.9cm,angle eccentricity=1.35]
      {angle=V--O--U};
  \draw[fill=black] (O) circle (1pt);
\end{tikzpicture}
\caption{Two arrows tail-to-tail at the origin enclose an angle $\theta$. The
third side of the triangle is $\vec u-\vec v$, and the law of cosines relates its
length to $\norm{\vec u}$, $\norm{\vec v}$, and $\cos\theta$.}
\label{ch04:loc}
\end{figure}

The law of cosines from trigonometry says that in a triangle with side lengths
$a$, $b$, $c$ and angle $\theta$ opposite the side of length $c$,
\[
  c^2 = a^2 + b^2 - 2ab\cos\theta.
\]
Apply it to our triangle, with $a = \norm{\vec u}$, $b = \norm{\vec v}$, and the
opposite side $c = \norm{\vec u - \vec v}$:
\[
  \norm{\vec u - \vec v}^2
    = \norm{\vec u}^2 + \norm{\vec v}^2 - 2\,\norm{\vec u}\,\norm{\vec v}\cos\theta.
\]
Now compute that left side with the dot product, using linearity and symmetry:
\begin{align*}
  \norm{\vec u - \vec v}^2
    &= (\vec u - \vec v)\cdot(\vec u - \vec v) \\
    &= \vec u\cdot\vec u - \vec u\cdot\vec v - \vec v\cdot\vec u + \vec v\cdot\vec v \\
    &= \norm{\vec u}^2 - 2\,(\vec u\cdot\vec v) + \norm{\vec v}^2.
\end{align*}
Set the two expressions for $\norm{\vec u-\vec v}^2$ equal. The
$\norm{\vec u}^2$ and $\norm{\vec v}^2$ on both sides cancel, leaving
\[
  -2\,(\vec u\cdot\vec v) = -2\,\norm{\vec u}\,\norm{\vec v}\cos\theta,
\]
and dividing by $-2$ gives precisely
\[
  \vec u\cdot\vec v = \norm{\vec u}\,\norm{\vec v}\,\cos\theta.
\]
That is the whole derivation. The algebraic dot product and the geometric ``length
times length times cosine'' are two faces of the same number, glued together by the
law of cosines.

\begin{aestheticbox}
Pause on how surprising this is. The left-hand side, $\sum_i u_iv_i$, is pure
arithmetic -- it never mentions angles. The right-hand side is pure geometry --
lengths and a cosine. The identity says these always agree. This is the bridge
between coordinates and geometry, and it is exactly the bridge that lets a computer,
which only ever knows lists of numbers, reason about angles, similarity, and
perpendicularity it can never ``see.''
\end{aestheticbox}

\subsection{Solving for the angle}

Turn the boxed formula around to get the angle from the coordinates. As long as
neither vector is zero,
\[
  \cos\theta = \frac{\vec u\cdot\vec v}{\norm{\vec u}\,\norm{\vec v}},
  \qquad
  \theta = \arccos\!\left(\frac{\vec u\cdot\vec v}
                                {\norm{\vec u}\,\norm{\vec v}}\right).
\]

\begin{example}\label{ch04:angle-ex}
Find the angle between $\vec u = \colvec{3\\4}$ and $\vec v = \colvec{2\\-1}$
(the vectors from Example~\ref{ch04:first-dot}). We already found
$\vec u\cdot\vec v = 2$. The lengths are
\[
  \norm{\vec u} = \sqrt{9+16} = 5, \qquad
  \norm{\vec v} = \sqrt{4+1} = \sqrt5 \approx 2.236.
\]
So
\[
  \cos\theta = \frac{2}{5\sqrt5} \approx 0.1789,
  \qquad
  \theta \approx \arccos(0.1789) \approx 79.7^\circ.
\]
The dot product was positive, so we expected an acute angle, and $79.7^\circ$
just barely qualifies -- the arrows are nearly, but not quite, perpendicular.
\end{example}

\begin{labbox}
Compute by hand, then sketch each pair as arrows and check the sign of the angle
against the sign of the dot product.
\begin{enumerate}
  \item $\colvec{1\\0}\cdot\colvec{1\\1}$ and the angle between them.
        (Answer: dot $=1$, angle $=45^\circ$.)
  \item $\colvec{2\\1}\cdot\colvec{-1\\2}$. What is the angle? What does the
        value of the dot product tell you before you compute the angle?
  \item $\colvec{1\\1\\0}\cdot\colvec{1\\1\\1}$, then $\cos\theta$, then $\theta$.
        (This is the angle between a face diagonal and the space diagonal of a
        cube.)
\end{enumerate}
\end{labbox}

\section{Orthogonality: perpendicular means dot product zero}

The cosine formula hands us the cleanest definition in linear algebra. Two nonzero
arrows are perpendicular exactly when the angle between them is $90^\circ$, i.e.
$\cos\theta = 0$. Since $\vec u\cdot\vec v = \norm{\vec u}\norm{\vec v}\cos\theta$
and the lengths are nonzero, the dot product is zero precisely when $\cos\theta = 0$:
\[
  \vec u \perp \vec v \quad\Longleftrightarrow\quad \vec u\cdot\vec v = 0.
\]

\begin{definition}\label{ch04:orthogonal}
Two vectors $\vec u, \vec v\in\R^n$ are \textbf{orthogonal} (the fancy word for
perpendicular) when $\vec u\cdot\vec v = 0$. By this definition the zero vector is
orthogonal to \emph{every} vector, since $\vec 0\cdot\vec v = 0$ always.
\end{definition}

We trade the geometric word ``perpendicular'' for the algebraic test ``dot product
$=0$,'' and the trade is a fantastic deal: the algebraic test works in any number
of dimensions, where we can no longer draw the right angle but the equation still
makes perfect sense. Two vectors in $\R^{100}$ can be orthogonal; you will never
picture it, but you can always check it.

\begin{figure}[H]
\centering
\begin{tikzpicture}[scale=1.1,>=Latex]
  \draw[->,gray!50] (-2.5,0) -- (2.8,0);
  \draw[->,gray!50] (0,-1) -- (0,3);
  \draw[->,very thick,vecblue] (0,0) -- (2,1)
        node[right]{$\vec a=\colvec{2\\1}$};
  \draw[->,very thick,vecred] (0,0) -- (-1,2)
        node[above left]{$\vec b=\colvec{-1\\2}$};
  % right-angle square at the origin, built from the two unit directions
  \draw[gray] (0.358,0.179) -- (0.179,0.537) -- (-0.179,0.358);
  \node[gray] at (0.42,0.46) {\small$90^\circ$};
  \draw[fill=black] (0,0) circle (1pt);
\end{tikzpicture}
\caption{$\vec a\cdot\vec b = (2)(-1)+(1)(2) = 0$, so the arrows meet at a right
angle. Orthogonality is just a zero dot product.}
\label{ch04:perp}
\end{figure}

\begin{example}
Back in Example~\ref{ch04:first-dot} we found
$\colvec{1\\2\\2}\cdot\colvec{2\\0\\-1}=0$. Those two three-dimensional arrows are
orthogonal -- they meet at a perfect right angle in space, even though it would take
a careful drawing to convince your eye. The algebra knew instantly.
\end{example}

\begin{strangbox}
\emph{Before Strang:} when he opens his lectures by saying ``perpendicular means
the dot product is zero,'' you will already own that sentence -- you will have
\emph{derived} it from the cosine formula rather than been told to memorize it.
\emph{After Strang:} this one idea, ``$\vec u\cdot\vec v=0$,'' is the seed of
orthogonal subspaces, the four fundamental subspaces meeting at right angles,
orthogonal matrices, projections, and the Gram--Schmidt process. Whenever Strang
draws two subspaces crossing at $90^\circ$, the dot product of this chapter is doing
the actual work underneath.
\end{strangbox}

\subsection{The Pythagorean theorem, reborn}

Orthogonality gives the cleanest possible version of Pythagoras. If
$\vec u\cdot\vec v = 0$, expand the squared length of the sum:
\[
  \norm{\vec u + \vec v}^2
    = (\vec u+\vec v)\cdot(\vec u+\vec v)
    = \norm{\vec u}^2 + 2\,(\vec u\cdot\vec v) + \norm{\vec v}^2.
\]
The middle term $2(\vec u\cdot\vec v)$ vanishes precisely when the vectors are
orthogonal, leaving
\[
  \norm{\vec u + \vec v}^2 = \norm{\vec u}^2 + \norm{\vec v}^2
  \qquad\text{exactly when } \vec u\perp\vec v.
\]
This is the Pythagorean theorem, stated for vectors, derived in two lines, and
valid in every dimension. The converse holds too: if
$\norm{\vec u+\vec v}^2 = \norm{\vec u}^2+\norm{\vec v}^2$, then the cross term must
be zero, so $\vec u\cdot\vec v = 0$ and the vectors are orthogonal. Pythagoras and
orthogonality are the \emph{same} statement.

\section{Unit vectors and direction}

Often we care only about the \emph{direction} an arrow points, not its length. To
strip out the length, divide a nonzero vector by its own norm:
\[
  \hat{\vec u} \;=\; \frac{\vec u}{\norm{\vec u}}.
\]
This is called \textbf{normalizing} $\vec u$, and the result $\hat{\vec u}$ (read
``$\vec u$ hat'') is a unit vector pointing the same way as $\vec u$, because
\[
  \norm{\hat{\vec u}} = \norm[\Big]{\frac{1}{\norm{\vec u}}\vec u}
  = \frac{1}{\norm{\vec u}}\,\norm{\vec u} = 1,
\]
using the scaling rule $\norm{c\vec u}=\abs c\,\norm{\vec u}$ with the positive
scalar $c = 1/\norm{\vec u}$.

\begin{example}
Normalize $\vec u = \colvec{3\\4}$. Its length is $5$, so
$\hat{\vec u} = \tfrac15\colvec{3\\4} = \colvec{3/5\\4/5} = \colvec{0.6\\0.8}$,
and indeed $0.6^2 + 0.8^2 = 0.36+0.64 = 1$. \checkmark
\end{example}

Unit vectors make the cosine formula especially clean. If $\hat{\vec u}$ and
$\hat{\vec v}$ are both unit vectors, then
$\norm{\hat{\vec u}}=\norm{\hat{\vec v}}=1$ and the formula collapses to
\[
  \hat{\vec u}\cdot\hat{\vec v} = \cos\theta.
\]
For unit vectors, \emph{the dot product simply is the cosine of the angle between
them}. Hold onto that -- it is the entire idea behind cosine similarity at the end
of the chapter.

\begin{focusbox}
Two complementary readings of the dot product:
\begin{itemize}
  \item $\vec u\cdot\vec u = \norm{\vec u}^2$ (length squared), and
  \item $\hat{\vec u}\cdot\hat{\vec v} = \cos\theta$ (cosine of the angle), once you
        have normalized.
\end{itemize}
Length and angle, the two pieces of metric geometry, both come straight out of this
one operation.
\end{focusbox}

\section{A first look at projection}

Here is a preview of one of the most important uses of the dot product; a whole
later chapter is devoted to it, but the seed belongs here. Ask: \emph{how much
of $\vec u$ points along the direction of $\vec v$?}

Picture shining a light straight down onto the line through $\vec v$ and reading off
the shadow $\vec u$ casts on that line (Figure~\ref{ch04:proj}). That shadow is the
\textbf{projection} of $\vec u$ onto $\vec v$.

\begin{figure}[H]
\centering
\begin{tikzpicture}[scale=1.4,>=Latex]
  \coordinate (O) at (0,0);
  \coordinate (V) at (3,0);
  \coordinate (U) at (1.8,1.5);
  \coordinate (P) at (1.8,0);
  \draw[->,gray!50] (-0.4,0) -- (3.3,0);
  \draw[->,very thick,vecred] (O) -- (V) node[below right]{$\vec v$};
  \draw[->,very thick,vecblue] (O) -- (U) node[above]{$\vec u$};
  \draw[->,thick,targetgreen] (O) -- (P)
        node[midway,below]{$\proj_{\vec v}\vec u$};
  \draw[dashed,gray] (U) -- (P);
  \draw (1.65,0) -- (1.65,0.15) -- (1.8,0.15);
  \pic[draw,"$\theta$",angle radius=0.6cm]{angle=V--O--U};
  \draw[fill=black] (O) circle (1pt);
\end{tikzpicture}
\caption{The projection of $\vec u$ onto $\vec v$ is the shadow $\vec u$ casts on the
line through $\vec v$. Its signed length is $\norm{\vec u}\cos\theta$.}
\label{ch04:proj}
\end{figure}

From the right triangle in the picture, the (signed) length of that shadow is the
adjacent side, $\norm{\vec u}\cos\theta$. But
$\norm{\vec u}\cos\theta = \dfrac{\vec u\cdot\vec v}{\norm{\vec v}}$ straight from
the cosine formula. So the length of the shadow is
\[
  \norm{\vec u}\cos\theta = \frac{\vec u\cdot\vec v}{\norm{\vec v}},
\]
and the shadow as an actual vector -- this length, pointed along the unit vector
$\hat{\vec v} = \vec v/\norm{\vec v}$ -- is
\[
  \proj_{\vec v}\vec u
   \;=\; \left(\frac{\vec u\cdot\vec v}{\norm{\vec v}}\right)\hat{\vec v}
   \;=\; \frac{\vec u\cdot\vec v}{\vec v\cdot\vec v}\,\vec v.
\]

\begin{example}
Project $\vec u = \colvec{1\\3}$ onto $\vec v = \colvec{2\\0}$. Here
$\vec u\cdot\vec v = 2$ and $\vec v\cdot\vec v = 4$, so
\[
  \proj_{\vec v}\vec u = \tfrac24\colvec{2\\0} = \tfrac12\colvec{2\\0}
                       = \colvec{1\\0}.
\]
That is exactly right: the shadow of $\colvec{1\\3}$ on the horizontal axis is
$\colvec{1\\0}$ -- it keeps the horizontal part and drops the vertical part, just as
the picture promised.
\end{example}

\begin{intuitionbox}
Every appearance of the fraction $\dfrac{\vec u\cdot\vec v}{\vec v\cdot\vec v}$ later
in this book -- in least squares, in regression, in Gram--Schmidt -- is this same
shadow. The dot product on top measures the overlap of $\vec u$ with $\vec v$; the
dot product on the bottom rescales by the length of $\vec v$ so that the answer does
not change if you happen to write $\vec v$ longer or shorter. Projection is the dot
product earning its keep.
\end{intuitionbox}

\section{Two great inequalities}

We have been writing
$\cos\theta = \dfrac{\vec u\cdot\vec v}{\norm{\vec u}\norm{\vec v}}$ as if it
obviously makes sense. But a cosine must live between $-1$ and $1$. For the formula
to hold, the fraction on the right had better never escape $[-1,1]$ -- that is,
we need $\abs{\vec u\cdot\vec v} \le \norm{\vec u}\norm{\vec v}$. This is the
celebrated \textbf{Cauchy--Schwarz inequality}, and remarkably it holds in every
dimension, with pure algebra, no pictures required.

\subsection{Cauchy--Schwarz}

\begin{theorem}[Cauchy--Schwarz]\label{ch04:cs}
For all $\vec u,\vec v\in\R^n$,
\[
  \abs{\vec u\cdot\vec v} \;\le\; \norm{\vec u}\,\norm{\vec v}.
\]
Equality holds if and only if $\vec u$ and $\vec v$ are parallel (one is a scalar
multiple of the other, including the case where one is $\vec 0$).
\end{theorem}

\begin{proof}
If $\vec v = \vec 0$ both sides are zero and there is nothing to prove, so assume
$\vec v\ne\vec 0$. The trick is to use positivity (property 3): the squared length of
\emph{any} vector is $\ge 0$. Consider, for a real number $t$, the vector
$\vec u - t\,\vec v$, and expand its squared length:
\[
  0 \;\le\; \norm{\vec u - t\vec v}^2
    = (\vec u - t\vec v)\cdot(\vec u - t\vec v)
    = \norm{\vec u}^2 - 2t\,(\vec u\cdot\vec v) + t^2\,\norm{\vec v}^2.
\]
Read the right-hand side as a quadratic in $t$, opening upward (the $t^2$ coefficient
$\norm{\vec v}^2$ is positive). It is $\ge 0$ for \emph{every} $t$, so it can have at
most one real root, which forces its discriminant to be $\le 0$:
\[
  \bigl(-2(\vec u\cdot\vec v)\bigr)^2 - 4\,\norm{\vec v}^2\,\norm{\vec u}^2 \le 0.
\]
That is $4(\vec u\cdot\vec v)^2 \le 4\,\norm{\vec u}^2\norm{\vec v}^2$. Divide by $4$
and take square roots:
\[
  \abs{\vec u\cdot\vec v} \le \norm{\vec u}\,\norm{\vec v}.
\]
Equality means the discriminant is exactly zero, i.e. the quadratic has a real root
$t_0$ with $\norm{\vec u - t_0\vec v}^2 = 0$, which (positivity again) means
$\vec u = t_0\vec v$: the vectors are parallel.
\end{proof}

\begin{aestheticbox}
That proof used \emph{only} the three properties in our Formal box -- symmetry,
linearity, positivity. It never once peeked at coordinates or drew an angle. So the
inequality holds not just in $\R^n$ but in \emph{any} setting with a
dot-product-like operation obeying those three rules: spaces of functions, spaces of
random variables, infinite-dimensional spaces. The same five lines prove that the
correlation coefficient in statistics lies in $[-1,1]$ and that overlap integrals in
quantum mechanics are bounded. One little quadratic, enormous reach.
\end{aestheticbox}

\begin{historybox}
The inequality carries several names. Augustin-Louis Cauchy proved the version for
finite sums of real numbers in 1821; Viktor Bunyakovsky extended it to integrals in
1859; and Hermann Schwarz rediscovered the integral form in 1888 in his work on
minimal surfaces. In much of the world it is therefore called the
Cauchy--Bunyakovsky--Schwarz inequality. What unifies all three versions is exactly
the abstraction in the box above: any operation with symmetry, linearity, and
positivity obeys it, whether it sums coordinates, integrates a product of functions,
or averages a product of random variables.
\end{historybox}

With Cauchy--Schwarz in hand, the angle formula is finally legitimate: the fraction
$\frac{\vec u\cdot\vec v}{\norm{\vec u}\norm{\vec v}}$ genuinely lies in $[-1,1]$, so
there really is an angle $\theta\in[0,\pi]$ with that cosine. \emph{This is why we
are allowed to define angle by the dot product even in dimensions we cannot picture.}

\subsection{The triangle inequality}

The second great inequality is the one your intuition already trusts: the straight
path is the shortest, so going from a tail to a tip directly is no longer than
detouring. In vector language,
$\norm{\vec u + \vec v} \le \norm{\vec u} + \norm{\vec v}$.

\begin{figure}[H]
\centering
\begin{tikzpicture}[scale=1.0,>=Latex]
  \coordinate (O) at (0,0);
  \coordinate (A) at (3.5,0.6);
  \coordinate (B) at (5,2.4);
  \draw[->,very thick,vecblue] (O) -- (A) node[midway,below]{$\vec u$};
  \draw[->,very thick,vecred]  (A) -- (B) node[midway,below right]{$\vec v$};
  \draw[->,thick,structurepurple] (O) -- (B)
        node[midway,above left]{$\vec u+\vec v$};
  \draw[fill=black] (O) circle (1.4pt);
\end{tikzpicture}
\caption{The direct trip $\vec u+\vec v$ is never longer than the detour
$\norm{\vec u}+\norm{\vec v}$ through the corner. That is the triangle inequality.}
\label{ch04:triangle}
\end{figure}

\begin{theorem}[Triangle inequality]\label{ch04:tri}
For all $\vec u,\vec v\in\R^n$,
\[
  \norm{\vec u + \vec v} \;\le\; \norm{\vec u} + \norm{\vec v}.
\]
\end{theorem}

\begin{proof}
Square the left side and expand:
\[
  \norm{\vec u+\vec v}^2
   = \norm{\vec u}^2 + 2\,(\vec u\cdot\vec v) + \norm{\vec v}^2.
\]
Now bound the cross term. Since $\vec u\cdot\vec v \le \abs{\vec u\cdot\vec v}$, and
by Cauchy--Schwarz $\abs{\vec u\cdot\vec v}\le\norm{\vec u}\norm{\vec v}$,
\[
  \norm{\vec u+\vec v}^2
   \le \norm{\vec u}^2 + 2\,\norm{\vec u}\norm{\vec v} + \norm{\vec v}^2
   = \bigl(\norm{\vec u} + \norm{\vec v}\bigr)^2.
\]
Both sides are nonnegative, so taking square roots preserves the inequality:
$\norm{\vec u+\vec v}\le\norm{\vec u}+\norm{\vec v}$.
\end{proof}

Notice how the triangle inequality \emph{rode in on the back of} Cauchy--Schwarz --
the cross term is exactly the dot product, and Cauchy--Schwarz is precisely the tool
that bounds it. The two inequalities are not separate facts; the second is a
corollary of the first.

\begin{walkbox}
Take a walk and turn this over: equality in the triangle inequality,
$\norm{\vec u+\vec v}=\norm{\vec u}+\norm{\vec v}$, means the detour is no longer than
the direct route -- the ``corner'' has flattened out. Geometrically, when must that
happen? Trace the proof: equality forces equality in Cauchy--Schwarz, which forces
$\vec u$ and $\vec v$ parallel, \emph{and} forces $\vec u\cdot\vec v\ge 0$, so they
point the \emph{same} way, not opposite ways. Does that match your picture of when
two arrows lie perfectly end to end along a single line?
\end{walkbox}

\begin{connectionbox}
The two properties $\norm{\vec u}\ge 0$ (with equality only at $\vec 0$) and
$\norm{\vec u+\vec v}\le\norm{\vec u}+\norm{\vec v}$, together with
$\norm{c\vec u}=\abs c\norm{\vec u}$, are exactly the axioms a mathematician demands
of any ``length'' -- they define what is called a \emph{norm}, and a norm in turn
defines a notion of \emph{distance}, $d(\vec u,\vec v) = \norm{\vec u-\vec v}$,
satisfying $d(\vec u,\vec w)\le d(\vec u,\vec v)+d(\vec v,\vec w)$. This is the
doorway from linear algebra into \emph{metric spaces, analysis, and topology}: the
moment you can measure distance and trust the triangle inequality, you can talk about
limits, continuity, and convergence. Our humble dot product is where all of that
begins.
\end{connectionbox}

\section{Modern application: cosine similarity}

Why does any of this matter to a search engine, a recommender, or a large language
model? Because those systems turn things -- documents, images, words -- into vectors,
and then must decide which vectors are \emph{similar}. The dot product, normalized,
is the standard measure of that similarity, and it works for exactly the reason we
just proved: the normalized dot product is a cosine, so it lives in $[-1,1]$ and
reads as ``how much do these point the same way.''

\begin{applicationbox}
\textbf{Cosine similarity.} Represent a document as a vector of word counts. Suppose
our vocabulary is just \{\texttt{linear}, \texttt{algebra}, \texttt{cooking}\} and
\[
  \vec d_1 = \colvec{4\\3\\0}\;(\text{a math note}),\quad
  \vec d_2 = \colvec{8\\6\\0}\;(\text{a longer math note}),\quad
  \vec d_3 = \colvec{0\\1\\5}\;(\text{a recipe}).
\]
The \textbf{cosine similarity} of two documents is the cosine of the angle between
their vectors,
\[
  \cos\theta = \frac{\vec d_i\cdot\vec d_j}{\norm{\vec d_i}\,\norm{\vec d_j}}.
\]
Compare the two math notes:
\[
  \vec d_1\cdot\vec d_2 = 32 + 18 + 0 = 50, \quad
  \norm{\vec d_1} = 5,\quad \norm{\vec d_2}=10,\quad
  \cos\theta = \frac{50}{50} = 1.
\]
A perfect score of $1$: even though $\vec d_2$ has twice the word counts (it is
simply a longer document), it points in exactly the same direction, so cosine
similarity declares the two notes identical in \emph{topic}. That
length-insensitivity is the whole point -- we do not want a document to look
different just because it is longer. Now compare the math note to the recipe:
\[
  \vec d_1\cdot\vec d_3 = 0 + 3 + 0 = 3,\quad
  \norm{\vec d_3} = \sqrt{26}\approx 5.10,\quad
  \cos\theta = \frac{3}{5\cdot 5.10}\approx 0.118.
\]
Nearly orthogonal -- a similarity near zero -- which is exactly what we want: the
recipe is about something else. The same arithmetic, run over vectors with hundreds
or thousands of coordinates (one per word, or per learned ``embedding'' dimension),
is how modern search, recommendation, and retrieval systems find the items closest in
meaning to your query. Underneath the buzzwords, it is $\sum_i u_iv_i$ divided by two
lengths.
\end{applicationbox}

To see why length-insensitivity is the right behaviour, drop the (here empty)
\texttt{cooking} axis and plot the three document vectors in the
\texttt{linear}--\texttt{algebra} plane (Figure~\ref{ch04:cosine}). The two math
notes $\vec d_1$ and $\vec d_2$ lie on the \emph{same ray} from the origin --
$\vec d_2$ is just $\vec d_1$ stretched -- so the angle between them is $0$ and the
cosine is $1$. The recipe $\vec d_3$ has almost all of its weight on the
\texttt{cooking} axis (out of the page), so once that axis is restored the full
angle to either math note is close to $90^\circ$ and the cosine is close to $0$.

\begin{figure}[H]
\centering
\begin{tikzpicture}[scale=0.55,>=Latex]
  \draw[->,gray!50] (-0.4,0) -- (9,0) node[right]{\small\texttt{linear} count};
  \draw[->,gray!50] (0,-0.4) -- (0,7) node[above]{\small\texttt{algebra} count};
  \draw[->,very thick,vecblue] (0,0) -- (4,3)
        node[above]{$\vec d_1$};
  \draw[->,very thick,vecblue!55] (0,0) -- (8,6)
        node[right]{$\vec d_2$};
  \draw[->,very thick,vecred] (0,0) -- (0,1)
        node[left]{$\vec d_3$};
  \node[align=left,anchor=west] at (4.6,2.1)
    {\small same direction:\\ \small $\cos\theta=1$};
  \draw[fill=black] (0,0) circle (1.5pt);
\end{tikzpicture}
\caption{The three document vectors projected onto the
\texttt{linear}--\texttt{algebra} plane (the \texttt{cooking} coordinate is the
out-of-page third axis, where only $\vec d_3$ has weight). Cosine similarity reads
\emph{direction}, not size: $\vec d_1$ and $\vec d_2$ share a ray and score $1$,
while $\vec d_3$ lives mostly along the out-of-page \texttt{cooking} axis, so in
full space it is nearly perpendicular to both and scores near $0$.}
\label{ch04:cosine}
\end{figure}

\begin{trapbox}
Cosine similarity ignores length \emph{on purpose}. Sometimes that is wrong for your
problem -- if magnitude carries real information (total sales, signal energy, how
many times a user clicked), then normalizing it away throws out the very thing you
cared about. Choose deliberately: cosine similarity (angle only) versus a raw dot
product or Euclidean distance (which feel length). The dot product gives you both
ingredients; \emph{you} decide which to keep.
\end{trapbox}

\section{Worked summary example}

Let us run one example through every idea in the chapter, end to end. Take
\[
  \vec u = \colvec{2\\2\\1}, \qquad \vec v = \colvec{0\\3\\4}.
\]

\textbf{Dot product.}
$\vec u\cdot\vec v = (2)(0)+(2)(3)+(1)(4) = 0+6+4 = 10.$

\textbf{Lengths.}
$\norm{\vec u} = \sqrt{4+4+1} = \sqrt9 = 3$, \quad
$\norm{\vec v} = \sqrt{0+9+16} = \sqrt{25} = 5.$

\textbf{Angle.}
$\cos\theta = \dfrac{10}{3\cdot5} = \dfrac{10}{15} = \dfrac23 \approx 0.667$,
so $\theta = \arccos(0.667)\approx 48.2^\circ$. Positive dot product, acute angle:
consistent.

\textbf{Cauchy--Schwarz check.}
$\abs{\vec u\cdot\vec v} = 10$ and $\norm{\vec u}\norm{\vec v} = 15$, and indeed
$10 \le 15$. \checkmark

\textbf{Triangle inequality check.}
$\vec u+\vec v = \colvec{2\\5\\5}$, so $\norm{\vec u+\vec v} = \sqrt{4+25+25}
= \sqrt{54}\approx 7.35$, while $\norm{\vec u}+\norm{\vec v} = 3+5 = 8$. Indeed
$7.35 \le 8$. \checkmark

\textbf{Unit vector.}
$\hat{\vec u} = \tfrac13\colvec{2\\2\\1} = \colvec{2/3\\2/3\\1/3}$, and
$(2/3)^2+(2/3)^2+(1/3)^2 = \tfrac{4+4+1}{9} = 1.$ \checkmark

\textbf{Projection of $\vec u$ onto $\vec v$.}
\[
  \proj_{\vec v}\vec u
   = \frac{\vec u\cdot\vec v}{\vec v\cdot\vec v}\,\vec v
   = \frac{10}{25}\colvec{0\\3\\4}
   = \tfrac25\colvec{0\\3\\4}
   = \colvec{0\\6/5\\8/5}.
\]
Sanity check: the shadow should be shorter than $\vec u$. Its length is
$\tfrac25\norm{\vec v} = \tfrac25\cdot 5 = 2$, which is indeed less than
$\norm{\vec u}=3$. \checkmark Every tool in the chapter, one vector pair, all
consistent.

\begin{recoverybox}
If that last example felt like a lot, that is the right feeling -- it deliberately
touched everything. The thing to carry forward is small: \textbf{one operation},
$\vec u\cdot\vec v = \sum_i u_iv_i$. From it, length is $\sqrt{\vec u\cdot\vec u}$,
angle is $\arccos\!\frac{\vec u\cdot\vec v}{\norm{\vec u}\norm{\vec v}}$,
perpendicular is $\vec u\cdot\vec v = 0$, and projection is
$\frac{\vec u\cdot\vec v}{\vec v\cdot\vec v}\vec v$. Four ideas, one engine.

\textbf{Cost.} A dot product in $\R^n$ is $n$ multiplications and $n-1$ additions --
linear in the dimension, $O(n)$ work. Length costs one dot product plus a square
root; an angle costs two more lengths and an $\arccos$; a projection costs two dot
products and a scaling. Everything in this chapter is cheap, which is exactly why
cosine similarity scales to vectors with thousands of coordinates and to millions of
documents.
\end{recoverybox}

\section{Exercises}

\begin{exercise}\label{ch04:ex-warmup}
Compute each dot product and say whether the angle between the vectors is acute,
right, or obtuse.
\[
  \text{(a) } \colvec{1\\2}\cdot\colvec{3\\1},\qquad
  \text{(b) } \colvec{2\\-3}\cdot\colvec{3\\2},\qquad
  \text{(c) } \colvec{1\\1\\1}\cdot\colvec{1\\1\\-3}.
\]
\end{exercise}

\begin{exercise}
Find the length of each vector and then normalize it to a unit vector.
\[
  \text{(a) } \colvec{6\\8}, \qquad
  \text{(b) } \colvec{1\\-2\\2}, \qquad
  \text{(c) } \colvec{1\\1\\1\\1}.
\]
\end{exercise}

\begin{exercise}
Find the angle (in degrees, to one decimal place) between
$\vec u=\colvec{1\\0\\1}$ and $\vec v=\colvec{1\\1\\0}$. These are two face diagonals
of a cube meeting at a corner; does your answer match the $60^\circ$ you might guess
from the geometry?
\end{exercise}

\begin{exercise}
For what value(s) of $k$ are $\colvec{2\\k}$ and $\colvec{3\\-1}$ orthogonal? For
what value(s) of $k$ are they parallel?
\end{exercise}

\begin{exercise}
Project $\vec u=\colvec{3\\4}$ onto $\vec v=\colvec{1\\0}$, and then onto
$\vec w=\colvec{0\\1}$. Add your two projection vectors. What do you get, and why
does it make sense?
\end{exercise}

\begin{exercise}
Let $\vec u=\colvec{2\\1}$ and $\vec v=\colvec{1\\3}$. Verify the Cauchy--Schwarz
inequality and the triangle inequality numerically for this pair. Are either of them
close to equality? Explain in one sentence what would have to be true for equality.
\end{exercise}

\begin{exercise}
Show that for any vectors $\vec u,\vec v\in\R^n$,
\[
  \norm{\vec u+\vec v}^2 + \norm{\vec u-\vec v}^2 = 2\norm{\vec u}^2 + 2\norm{\vec v}^2.
\]
(This is the \emph{parallelogram law}.) Then draw a parallelogram with sides
$\vec u$ and $\vec v$ and explain which lengths the four terms represent.
\emph{Hint:} expand both squared norms using the dot product and watch the cross
terms.
\end{exercise}

\begin{exercise}
Prove that if $\vec u$ is orthogonal to both $\vec v$ and $\vec w$, then $\vec u$ is
orthogonal to every linear combination $a\vec v + b\vec w$. \emph{Hint:} use
linearity of the dot product. (This little fact is the reason orthogonality plays so
nicely with subspaces later on.)
\end{exercise}

\begin{exercise}
A unit vector in the plane can be written
$\hat{\vec u}=\colvec{\cos\alpha\\\sin\alpha}$ and another as
$\hat{\vec v}=\colvec{\cos\beta\\\sin\beta}$. Compute
$\hat{\vec u}\cdot\hat{\vec v}$ directly from the coordinates, and use the fact that
this equals $\cos\theta$ (with $\theta$ the angle between them) to re-derive the
angle-subtraction formula
$\cos(\alpha-\beta)=\cos\alpha\cos\beta+\sin\alpha\sin\beta$. A trig identity,
falling out of the dot product.
\end{exercise}

\begin{exercise}
Three documents over the vocabulary \{\texttt{cat}, \texttt{dog}, \texttt{stock},
\texttt{market}\} are
\[
  \vec a=\colvec{3\\2\\0\\0},\quad
  \vec b=\colvec{6\\4\\0\\0},\quad
  \vec c=\colvec{0\\0\\4\\3}.
\]
Compute the cosine similarity of each pair. Which two documents are most similar, and
why does $\vec a$ versus $\vec b$ score $1$ even though their word counts differ?
\end{exercise}

\begin{exercise}[Challenge]
Use Cauchy--Schwarz to prove that for any real numbers $a_1,\dots,a_n$,
\[
  (a_1 + a_2 + \cdots + a_n)^2 \;\le\; n\,(a_1^2 + a_2^2 + \cdots + a_n^2).
\]
\emph{Hint:} apply the inequality to $\vec a = \colvec{a_1\\\vdots\\a_n}$ and a very
well-chosen second vector. When is it an equality? (This is one route to the fact
that the average of squares is at least the square of the average.)
\end{exercise}

\begin{exercise}[Challenge]
The angle formula lets us define angle in any dimension. Let $\vec u$, $\vec v$,
$\vec w$ be three unit vectors and let
$\theta_{uv},\theta_{vw},\theta_{uw}$ be the pairwise angles. Argue, using the
triangle inequality applied to arcs on the unit sphere (or the cosine formula and a
little trig), that $\theta_{uw}\le\theta_{uv}+\theta_{vw}$. Why is this the
``triangle inequality for angles,'' and where have you seen its flat-space cousin in
this chapter?
\end{exercise}
